\clearpage
\section{15分钟自测：题目与完整答案}
以下均为\textbf{补充自编练习}，用于检查掌握程度，不是老师原题，也不代表实际考题。每组建议先独立做15分钟，再核对答案。
\subsection{自测A：分辨率与滤波选择}
\subsubsection*{题1：相机与线扫描的需求}
生产线上待拍物体的横向视场为$1.2\,\mathrm{m}$，希望检测宽度$0.3\,\mathrm{mm}$的缺陷。假定每个缺陷在横向至少占3个像素。物体以$0.8\,\mathrm{m/s}$运动，并要求纵向采样间距不大于横向采样间距。
\begin{enumerate}
\item 横向至少需要多少像素？
\item 线扫描相机至少需要多大的行频？
\item 如果最大允许运动模糊为1个像素，曝光时间上限是多少？
\item 为什么满足这些数值还不能保证缺陷一定能被检出？
\end{enumerate}
\subsubsection*{题2：如何处理噪声}
某灰度图像大部分区域平缓，但有少量随机的纯黑点和纯白点。有人建议直接用拉普拉斯锐化，另一个人建议先用$3\times3$中值滤波。解释哪种建议更合理，并说明中值滤波的风险。
\subsubsection*{自测A参考答案}
\textbf{题1。}统一单位为毫米。单像素对应的物理距离至多为
\[
\Delta x=\frac{0.3}{3}=0.1\ \mathrm{mm/pixel}.
\]
横向像素数至少为
\[
N_x\ge\left\lceil\frac{1200}{0.1}\right\rceil=12000.
\]
物体速度$v=800\,\mathrm{mm/s}$，一行到下一行的位移为$v/f_{\mathrm{line}}$，因此
\[
f_{\mathrm{line}}\ge\frac{800}{0.1}=8000\ \mathrm{lines/s}.
\]
曝光期间的位移不超过一个像素：
\[
vT_{\mathrm{exp}}\le\Delta x
\quad\Rightarrow\quad
T_{\mathrm{exp}}\le\frac{0.1}{800}\ \mathrm{s}=125\ \mu\mathrm{s}.
\]
这些只是采样与运动模糊方面的必要设计约束。还要考虑镜头分辨能力、对焦、照明、信噪比、缺陷对比度、传送带滑动、触发同步及算法性能。不能把“每个缺陷有3个像素”当成普适检测保证。

\textbf{题2。}这种现象符合椒盐噪声的典型外观。拉普拉斯对突变敏感，会同时放大噪点，直接锐化不合理。小邻域中值滤波通过排序选择中间值，对稀疏极端离群值更稳健。但它可能消除真实的小亮点、细线，改变边界位置；精密尺寸测量时要验证偏差。也不能只凭“黑白点”就断言所有噪声都符合独立椒盐模型。
\clearpage
\subsection{自测B：灰度映射与直方图}
\subsubsection*{题1：gamma变换}
输入已归一化到$[0,1]$，使用$s=r^{\gamma}$。求$r=0.04,0.25,0.81$在$\gamma=0.5$时的输出，说明暗部发生什么变化。如果$\gamma=2$，$r=0.25$会变成多少？
\subsubsection*{题2：离散直方图均衡化}
一幅图只有4个灰度级$0,1,2,3$，对应像素数为$4,2,1,1$。使用
\[
s_k=\operatorname{round}\left[3\sum_{j=0}^{k}p(r_j)\right]
\]
进行均衡化，规定$1.5$四舍五入为$2$。求映射及输出直方图，并解释为什么输出不一定均匀。
\subsubsection*{自测B参考答案}
\textbf{题1。}当$\gamma=0.5$时，开平方得到
\[
0.04\mapsto0.2,\qquad0.25\mapsto0.5,\qquad0.81\mapsto0.9.
\]
在$0<r<1$范围内输出大于输入，图像整体趋于变亮；靠近暗部的输入区间被拉开，亮部相对压缩。$\gamma=2$时，$0.25^2=0.0625$，趋于变暗。必须说明归一化与参数定义；若软件使用$r^{1/\gamma}$，对“gamma增大”的解释会反过来。

\textbf{题2。}总像素数$8$，概率为
\[
\left[\frac12,\frac14,\frac18,\frac18\right],
\quad\text{累积分布为}\quad
\left[\frac12,\frac34,\frac78,1\right].
\]
乘3得到$[1.5,2.25,2.625,3]$，按题定取整规则得到映射
\[
0\mapsto2,\quad1\mapsto2,\quad2\mapsto3,\quad3\mapsto3.
\]
所以输出灰度$0,1,2,3$的像素数为$[0,0,6,2]$。离散灰度、有限样本和多对一映射使输出不必严格均匀。题目没有采用减去最小非零CDF的变体，不应自行更换公式。两种实现的结果可能不同，必须先说清约定。
\subsection{频域补充自测：不能只背“低通去噪”}
\textbf{补充题。}图像含细小真实边缘和周期性条纹噪声。为什么全局强低通可能不是最好选择？如果频谱中发现一对与原点对称的异常峰，怎样处理更有针对性？

\textbf{参考答案。}强低通会压低高频噪声，也会损失真实边缘和纹理；若条纹频率较低，还可能无法充分去除。可以考虑围绕那对异常峰设计成对陷波抑制，尽量保留其他频率。实图像频谱有共轭对称性，滤波器应保持相应对称关系，使反变换结果保持实数。陷波不能过宽，否则仍会删除有用信息；要确认异常峰确来自噪声而非真实周期纹理。
\clearpage
\section{最后一遍：看到题目先问自己什么}
\begin{enumerate}
\item \textbf{单位统一了吗？}米、毫米、微米，秒、毫秒、微秒必须一致；像素数取足够大的整数。
\item \textbf{题目要的是采样还是量化？}前者关注空间位置的离散化，后者关注幅值的离散化。
\item \textbf{灰度是否归一化？}不要把$[0,255]$中的数直接代入仅对$[0,1]$约定的公式。
\item \textbf{映射曲线的斜率说明什么？}局部斜率大于1意味着该灰度区间被扩展，不等于所有像素都变亮。
\item \textbf{直方图丢掉了什么？}它没有保留像素的空间排列，不能仅凭直方图判断物体位置或图像是否锐利。
\item \textbf{滤波核是否需要翻转？}相关与卷积对对称核给同样结果，对非对称核不一定相同。
\item \textbf{核权重和是多少？}归一化平滑核对常量图像保持常量；零和微分核对常量响应为零。
\item \textbf{滤波器对噪声模型合适吗？}去噪、细节保留、边界位置与计算成本要同时考虑。
\item \textbf{频率高低指什么？}是空间灰度变化快慢，不是“高频就是更亮”。
\item \textbf{结果需要截断吗？}锐化可能产生负值或超过灰度上限；先算完整结果，再按规定显示或裁剪，避免无符号整数溢出。
\end{enumerate}
